\Needspace{6\baselineskip}
\begin{center}\hypertarget{native-249-article177}{}\pdfbookmark[2]{177.}{article-177}177.\end{center}
Hinc iam sponte sequitur, systema maxime probabile valorum quantitatum
$p$, $q$, $r$, $s$ etc. id fore, in quo $\Omega$ valorem maximum obtineat,
adeoque ex $\nu$ aequationibus $\frac{\mathrm{d}\Omega}{\mathrm{d}p}=0$,
$\frac{\mathrm{d}\Omega}{\mathrm{d}q}=0$,
$\frac{\mathrm{d}\Omega}{\mathrm{d}r}=0$,
$\frac{\mathrm{d}\Omega}{\mathrm{d}s}=0$ etc. eruendum esse. Hae aequationes,
statuendo $M-V=v$, $M'-V'=v'$, $M''-V''=v''$ etc., atque
$\frac{\mathrm{d}\varphi\Delta}{\varphi\Delta\,.\,\mathrm{d}\Delta}=\varphi'\Delta$,
formam sequentem nanciscuntur:
\[
\begin{gathered}
\frac{\mathrm{d}v}{\mathrm{d}p}\varphi'v+
\frac{\mathrm{d}v'}{\mathrm{d}p}\varphi'v'+
\frac{\mathrm{d}v''}{\mathrm{d}p}\varphi'v''+\text{etc.}=0\\
\frac{\mathrm{d}v}{\mathrm{d}q}\varphi'v+
\frac{\mathrm{d}v'}{\mathrm{d}q}\varphi'v'+
\frac{\mathrm{d}v''}{\mathrm{d}q}\varphi'v''+\text{etc.}=0\\
\frac{\mathrm{d}v}{\mathrm{d}r}\varphi'v+
\frac{\mathrm{d}v'}{\mathrm{d}r}\varphi'v'+
\frac{\mathrm{d}v''}{\mathrm{d}r}\varphi'v''+\text{etc.}=0\\
\frac{\mathrm{d}v}{\mathrm{d}s}\varphi'v+
\frac{\mathrm{d}v'}{\mathrm{d}s}\varphi'v'+
\frac{\mathrm{d}v''}{\mathrm{d}s}\varphi'v''+\text{etc.}=0\\
\text{etc.}
\end{gathered}
\]

Hinc itaque per eliminationem problematis solutio plene determinata derivari
poterit, quamprimum functionis $\varphi'$ indoles innotuit. Quae quoniam a
priori definiri nequit, rem ab altera parte aggredientes inquiremus, cuinam
\SourceLeaf{250}functioni, tacite quasi pro basi acceptae, proprie innixum sit
principium trivium, cuius praestantia generaliter agnoscitur. Axiomatis
scilicet loco haberi solet hypothesis, si quae quantitas per plures
observationes immediatas, sub aequalibus circumstantiis aequalique cura
institutas, determinata fuerit, medium arithmeticum inter omnes valores
observatos exhibere valorem maxime probabilem, si non absoluto rigore, tamen
proxime saltem, ita ut semper tutissimum sit illi inhaerere. Statuendo itaque
$V=V'=V''=\text{etc.}=p$, generaliter esse debebit
$\varphi'(M-p)+\varphi'(M'-p)+\varphi'(M''-p)+\text{etc.}=0$, si pro $p$
substituitur valor $\frac{1}{\mu}(M+M'+M''+\text{etc.})$, quemcunque integrum
positivum exprimat $\mu$. Supponendo itaque $M'=M''=\text{etc.}=M-\mu N$,
erit generaliter, i. e. pro quovis valore integro positivo ipsius $\mu$,
$\varphi'(\mu-1)N=(1-\mu)\varphi'(-N)$, unde facile colligitur, generaliter
esse debere $\frac{\varphi'\Delta}{\Delta}$ quantitatem constantem, quam per
$k$ designabimus. Hinc fit $\log\varphi\Delta=\tfrac12 k\Delta\Delta+\text{Const.}$,
sive designando basin logarithmorum hyperbolicorum per $e$, supponendoque
$\text{Const.}=\log\varkappa$,
\[
\varphi\Delta=\varkappa e^{\frac12 k\Delta\Delta}.
\]
Porro facile perspicitur, $k$ necessario negativam esse debere, quo $\Omega$
revera fieri possit maximum, quamobrem statuemus $\tfrac12 k=-hh$; et quum
per theorema elegans primo ab ill. \textsc{Laplace} inventum, integrale
$\int e^{-hh\Delta\Delta}\mathrm{d}\Delta$, a $\Delta=-\infty$ usque ad
$\Delta=+\infty$, fiat $=\frac{\BareRoot{\pi}}{h}\mathpunct{\raisebox{.45ex}{,}}$
(denotando per $\pi$ semicircumferentiam circuli cuius radius $1$), functio
nostra fiet
\[
\varphi\Delta=\frac{h}{\BareRoot{\pi}}e^{-hh\Delta\Delta}.
\]
